%% ============================================================================
%%  chapters/chapter3.tex
%%  فصل ۳: روش‌شناسی و مدل پیشنهادی
%%  قالب رسمی پایان‌نامه و رساله دانشگاه خوارزمی
%% ============================================================================

\chapter{روش‌شناسی و مدل پیشنهادی}
\label{chap:methodology}

\section{مقدمه}
\label{sec:ch3_intro}

دستیابی به یک ساختار حروف‌چینی پایدار و خودکار برای پایان‌نامه‌ها و رساله‌های دانشگاهی، نیازمند طراحی یک متدولوژی علمی مدون و معماری چندلایه‌ای است که بتواند پیچیدگی‌های مربوط به فونت‌های فارسی، قوانین چیدمان دوسویه، شماره‌گذاری سلسله‌مراتبی، و فرمول‌بندی‌های پیشرفته ریاضی را مدیریت کند. در این فصل، ابتدا معماری کلان سامانه پیشنهادی و تفکیک لایه‌های کارکردی آن تشریح می‌گردد. سپس مدل‌سازی و فرمول‌بندی ریاضی اجزای سامانه شامل توابع هدف، روابط بهینه‌سازی ماتریسی، توابع تکه‌ای و انتگرال‌های احتمالاتی تبیین می‌شود. در ادامه، الگوریتم جامع پردازش و حل بهینه‌سازی بر پایه ساختار الگوریتمی ارائه‌شده و مشخصات فنی پیاده‌سازی ماژولار قالب به تفصیل بیان خواهد شد.

\section{معماری کلی سامانه و ساختار ماژولار}
\label{sec:ch3_architecture}

معماری سامانه پیشنهادی بر پایه اصل تفکیک دغدغه‌ها\LTRfootnote{\lr{Separation of Concerns (SoC)}} بنا نهاده شده است. بر این اساس، ساختار کلی پایان‌نامه به ۵ لایه مجزا تقسیم می‌گردد:
\begin{enumerate}
  \item \textbf{لایه هسته ساختاری و سندپایه (\lr{\texttt{kharazmi-thesis.cls}}):} این لایه وظیفه بارگذاری کلاس پایه گزارش، اعمال دقیق ابعاد هندسی کاغذ \lr{A4}، تعریف حاشیه‌ها (بالا ۳/۵، پایین ۲/۵، راست ۳/۰ و چپ ۳/۰ سانتی‌متر)، فاصله سرصفحه‌ها و پاصفحه‌ها، قالب‌بندی سلسله‌مراتب عناوین فصل‌ها و بخش‌ها، و قالب‌بندی زیرنویس جدول‌ها و شکل‌ها را بر عهده دارد.
  \item \textbf{لایه استایل و فراداده (\lr{\texttt{kharazmi.sty}}):} در این لایه، تنظیمات قلم‌های رسمی فارسی (بی زر ۱۳ برای متن و بی تیتر برای سرعنوان‌ها)، مقیاس‌بندی دقیق قلم لاتین تایمز نیو رومن، ماکروهای نگهداری مشخصات پایان‌نامه، و توابع تولیدکننده برگه‌های رسمی پیش از متن و پس از متن تعریف شده‌اند.
  \item \textbf{لایه صفحات آغازین (\lr{\texttt{front/}}):} شامل ماژول‌های مجزای جلد فارسی، صفحه بسم‌الله، صورت‌جلسه دفاع و تصویب‌نامه هیئت داوران، تعهدنامه اصالت اثر، تقدیم، سپاسگزاری، و چکیده فارسی است.
  \item \textbf{لایه فصول محتوایی (\lr{\texttt{chapters/}}):} ۵ فصل مستقل پایان‌نامه که به صورت فایل‌های مجزا توسط دستور \lr{\texttt{\textbackslash include}} به فایل اصلی متصل می‌شوند.
  \item \textbf{لایه مراجع و صفحات پایانی (\lr{\texttt{back/}} و \lr{\texttt{references.bib}}):} دربرگیرنده مراجع دوزبانه، پیوست‌های ریاضی و اثبات‌ها، واژه‌نامه‌های سه‌گانه، کارنامک علمی، و برگه‌های انگلیسی چکیده و جلد پشتی است.
\end{enumerate}

\section{فرمول‌بندی ریاضی و مدل‌سازی تحلیلی}
\label{sec:ch3_formulation}

در این بخش، به منظور ارائه نمونه‌ای کامل از قابلیت‌های پیشرفته ریاضی قالب دانشگاه خوارزمی، انواع گوناگون روابط ریاضی شامل معادلات تک‌خطی، چندخطی هم‌تراز، ماتریس‌ها، توابع تکه‌ای، و انتگرال‌ها فرمول‌بندی می‌گردد. تمامی شماره‌گذاری‌های روابط با استفاده از شناسه فصل و شماره رابطه به صورت فارسی (مانند \eqref{eq:ch3_single}) حروف‌چینی می‌شوند.

\subsection{معادلات تک‌خطی و تابع هدف}
فرض کنید $\mathcal{D} = \{(x_i, y_i)\}_{i=1}^{N}$ مجموعه داده‌های آموزشی شامل $N$ نمونه مستقل باشد. تابع هدف بهینه‌سازی\index{بهینه‌سازی} پارامترهای مدل $\theta \in \mathbb{R}^d$ با در نظر گرفتن ترم تنظیم‌سازی\index{تنظیم‌سازی} تیخونوف\LTRfootnote{\lr{Tikhonov Regularization}} به صورت رابطه تک‌خطی \eqref{eq:ch3_single} تعریف می‌شود:
\begin{equation}
\label{eq:ch3_single}
\mathcal{J}(\theta) = \frac{1}{N} \sum_{i=1}^{N} \ell\left(f(x_i; \theta), y_i\right) + \frac{\lambda}{2} \|\theta\|_2^2
\end{equation}
که در آن $\ell(\cdot, \cdot)$ نشان‌دهنده تابع زیان\index{تابع زیان} محلی بر روی هر نمونه، $\|\theta\|_2 = \sqrt{\sum_{j=1}^d \theta_j^2}$ بیانگر نُرم اقلیدسی بردار پارامترها، و $\lambda > 0$ ضریب کنترل پیچیدگی مدل است.

\subsection{معادلات چندخطی هم‌تراز}
برای به‌روزرسانی پارامترها در فرآیند گرادیان کاهشی\index{گرادیان کاهشی} تصادفی با تکانه مومنتوم نسترُف\LTRfootnote{\lr{Nesterov Accelerated Gradient}}، دنباله محاسبات به صورت رابطه چندخطی هم‌تراز \eqref{eq:ch3_multiline} صورت می‌پذیرد:
\begin{align}
\label{eq:ch3_multiline}
v_{t+1} &= \mu \, v_t - \eta \, \nabla_{\theta} \mathcal{J}\left(\theta_t + \mu \, v_t\right) \\
\theta_{t+1} &= \theta_t + v_{t+1} \\
\Delta \theta_t &= \|\theta_{t+1} - \theta_t\|_2 = \|v_{t+1}\|_2
\end{align}
که در آن $\mu \in [0, 1)$ ضریب تکانه، $\eta > 0$ نرخ یادگیری، و $v_t$ بردار سرعت در تکرار $t$ است.

\subsection{مدل‌سازی ماتریسی و سازوکار توجه خودکار}
در مدل‌های مدرن مبتنی بر توجه خودکار، ماتریس‌های پرسش ($\mathbf{Q}$)، کلید ($\mathbf{K}$) و مقدار ($\mathbf{V}$) از ضرب ماتریس ویژگی‌های ورودی $\mathbf{X} \in \mathbb{R}^{n \times d}$ در ماتریس‌های وزن پدید می‌آیند:
\begin{equation}
\label{eq:ch3_matrices}
\begin{bmatrix}
\mathbf{Q} \\
\mathbf{K} \\
\mathbf{V}
\end{bmatrix}
=
\begin{bmatrix}
\mathbf{W}_Q \\
\mathbf{W}_K \\
\mathbf{W}_V
\end{bmatrix}
\mathbf{X}^\top
\end{equation}
خروجی لایه توجه مقیاس‌شده به صورت ضرب ماتریسی در رابطه \eqref{eq:ch3_attention} مدل‌سازی می‌گردد:
\begin{equation}
\label{eq:ch3_attention}
\operatorname{Attention}(\mathbf{Q}, \mathbf{K}, \mathbf{V}) = \operatorname{softmax}\left( \frac{\mathbf{Q} \mathbf{K}^\top}{\sqrt{d_k}} \right) \mathbf{V}
\end{equation}
به طوری که ماتریس هسته توجه $\mathbf{A} \in \mathbb{R}^{n \times n}$ دارای درایه‌هایی به صورت زیر است:
\begin{equation}
\label{eq:ch3_softmax_matrix}
A_{ij} = \frac{\exp\left(\frac{q_i^\top k_j}{\sqrt{d_k}}\right)}{\sum_{l=1}^n \exp\left(\frac{q_i^\top k_l}{\sqrt{d_k}}\right)}
\end{equation}

\subsection{توابع چندضابطه‌ای (تکه‌ای)}
توابع فعال‌ساز غیرخطی در لایه‌های پنهان شبکه‌های عصبی نقش اساسی دارند. تابع فعال‌ساز اصلاح‌شده تکه‌ای خطی\LTRfootnote{\lr{Parametric Rectified Linear Unit (PReLU)}} به صورت رابطه \eqref{eq:ch3_piecewise} تعریف می‌شود:
\begin{equation}
\label{eq:ch3_piecewise}
\phi(z) =
\begin{cases}
z & \text{اگر } z \ge 0 \\
\alpha z & \text{اگر } z < 0
\end{cases}
\end{equation}
همچنین تابع زیان هابر\LTRfootnote{\lr{Huber Loss Function}} که نسبت به داده‌های پرت مقاوم است، ضابطه‌ای دوشاخه دارد:
\begin{equation}
\label{eq:ch3_huber}
\mathcal{L}_\delta(e) =
\begin{cases}
\frac{1}{2} e^2 & \text{اگر } |e| \le \delta \\[1em]
\delta \left( |e| - \frac{1}{2} \delta \right) & \text{وگرنه}
\end{cases}
\end{equation}

\subsection{روابط انتگرالی و محاسبات پیوسته}
در تحلیل‌های احتمالاتی و برآورد درست‌نمایی بیزی، انتگرال‌های چندبعدی بر روی توزیع پیشین پارامترها به کار می‌روند. امید ریاضی تابع زیان روی فضای پیوسته به صورت رابطه \eqref{eq:ch3_integral} محاسبه می‌گردد:
\begin{equation}
\label{eq:ch3_integral}
\mathbb{E}_{x \sim p(x)}\left[\mathcal{L}(x)\right] = \int_{\mathcal{X}} \mathcal{L}(x) \, p(x) \, \mathrm{d}x = \int_{-\infty}^{+\infty} \cdots \int_{-\infty}^{+\infty} \mathcal{L}(x_1, \dots, x_d) \, p(x_1, \dots, x_d) \, \mathrm{d}x_1 \cdots \mathrm{d}x_d
\end{equation}
همچنین احتمال پسین به وسیله قاعده بیز به صورت زیر فرمول‌بندی می‌شود:
\begin{equation}
\label{eq:ch3_bayes}
p(\theta \mid \mathcal{D}) = \frac{p(\mathcal{D} \mid \theta) \, p(\theta)}{\int_{\Theta} p(\mathcal{D} \mid \theta') \, p(\theta') \, \mathrm{d}\theta'}
\end{equation}

\section{طراحی الگوریتم پیشنهادی}
\label{sec:ch3_algorithm}

الگوریتم جامع بهینه‌سازی ساختار سند و آموزش مدل پیشنهادی در الگوریتم \ref{alg:ch3_optimization} با کلمات کلیدی فارسی به تصویر کشیده شده است. این الگوریتم با دریافت داده‌ها و فراسنج‌های اولیه، در هر دور فرآیند گرادیان کاهشی را اعمال کرده و شرایط توقف زودهنگام را می‌آزماید.

\begin{algorithm}[htbp]
\caption{آموزش و بهینه‌سازی پارامترهای مدل پیشنهادی}
\label{alg:ch3_optimization}
\begin{algorithmic}[1]
\REQUIRE مجموعه داده آموزشی $\mathcal{D}$، نرخ یادگیری $\eta$، تعداد کل دورها $E$، آستانه همگرایی $\epsilon$
\ENSURE بردار بهینه پارامترها $\theta^*$ و تاریخچه تابع زیان $\mathcal{H}$
\STATE بردار پارامترها را به صورت تصادفی مقداردهی اولیه کن: $\theta \leftarrow \theta_0$
\STATE متغیر تاریخچه را خالی در نظر بگیر: $\mathcal{H} \leftarrow \emptyset$
\FOR{$e = 1$ \TO $E$}
  \STATE داده‌های آموزشی $\mathcal{D}$ را به صورت تصادفی بر بزن
  \FOR{هر زیردسته داده $\mathcal{B} \subset \mathcal{D}$ با اندازه $B$}
    \STATE گرادیان تابع زیان روی زیردسته را محاسبه کن: $g \leftarrow \frac{1}{B} \sum_{x \in \mathcal{B}} \nabla_\theta \ell(x; \theta)$
    \STATE پارامترها را به‌روزرسانی کن: $\theta \leftarrow \theta - \eta \cdot g$
  \ENDFOR
  \STATE مقدار کل زیان را در پایان دور محاسبه کن: $L_e \leftarrow \mathcal{J}(\theta)$
  \STATE مقدار جدید را به تاریخچه بیفزا: $\mathcal{H} \leftarrow \mathcal{H} \cup \{L_e\}$
  \IF{$e > 1$ \AND $|L_e - L_{e-1}| < \epsilon$}
    \PRINT پیام همگرایی پیش از موعد در دور $e$
    \RETURN $\theta$
  \ENDIF
\ENDFOR
\RETURN $\theta$
\end{algorithmic}
\end{algorithm}

\section{پیاده‌سازی ماژولار و رابط‌های نرم‌افزاری}
\label{sec:ch3_interfaces}

به منظور اطمینان از سازگاری و انعطاف‌پذیری قالب، اجزای مختلف از طریق واسط‌های استاندارد با یکدیگر تعامل دارند. هر بخش از مستند دارای پرونده مستقل است و کاربر بدون نیاز به باز کردن یا دستکاری کدهای بنیادین کلاس، تنها با ویرایش فراداده‌های موجود در پرونده \lr{\texttt{main.tex}} اطلاعات شخصی، اسامی اساتید، عنوان و تاریخ را مقداردهی می‌کند. ساختار کلاس به گونه‌ای محافظت شده است که در صورت افزودن بسته‌های تکمیلی، هماهنگی سرصفحه‌ها و حاشیه‌های مصوب دانشگاه خوارزمی مخدوش نشود.

\section{جمع‌بندی فصل}
\label{sec:ch3_summary}

در این فصل، معماری ۵ لایه سامانه حروف‌چینی اسناد دانشگاهی تشریح و فرمول‌بندی تحلیلی انواع روابط ریاضی شامل معادلات تک‌خطی، چندخطی هم‌تراز، ماتریس‌ها، توابع تکه‌ای، و انتگرال‌های پیوسته به همراه شماره‌گذاری استاندارد فارسی ارائه گردید. همچنین ساختار الگوریتمی آموزش و بهینه‌سازی با کلمات کلیدی بومی‌سازی‌شده تدوین شد. در فصل چهارم، نتایج حاصل از ارزیابی عملی این مدل، جداول آماری، نمودارها، شکل‌ها و قطعه کدها به تفصیل تحلیل خواهند شد.
